\documentclass[11pt,a4paper]{report}

% ============================================================
% PACKAGES
% ============================================================

\usepackage[utf8]{inputenc}
\usepackage[T1]{fontenc}
\usepackage[french]{babel}

\usepackage{lmodern}
\usepackage{microtype}

\usepackage{amsmath}
\usepackage{amssymb}

\usepackage{graphicx}
\usepackage{float}
\usepackage{subcaption}

\usepackage{tikz}
\usetikzlibrary{positioning, arrows.meta}

\usepackage{booktabs}
\usepackage{array}
\usepackage{siunitx}

\usepackage{xcolor}
\usepackage{setspace}

% ============================================================
% MARGES
% ============================================================

\usepackage[
    a4paper,
    left=2.0cm,
    right=2.0cm,
    top=2.0cm,
    bottom=2.0cm,
    headheight=15pt,
    headsep=0.4cm
]{geometry}

% ============================================================
% EN-TETE ET PIED DE PAGE
% ============================================================

\usepackage{fancyhdr}

\pagestyle{fancy}
\fancyhf{}

\fancyhead[L]{\small ESRF -- High-performance algorithms for tomography}
\fancyhead[R]{\small Rana Alouani}

\fancyfoot[C]{\thepage}

\renewcommand{\headrulewidth}{0.4pt}
\renewcommand{\footrulewidth}{0pt}

% Même en-tête sur les premières pages de chapitre
\fancypagestyle{plain}{
    \fancyhf{}

    \fancyhead[L]{\small ESRF -- High-performance algorithms for tomography}
    \fancyhead[R]{\small Rana Alouani}

    \fancyfoot[C]{\thepage}

    \renewcommand{\headrulewidth}{0.4pt}
    \renewcommand{\footrulewidth}{0pt}
}

% ============================================================
% TITRES
% ============================================================

\usepackage{titlesec}

\titleformat{\chapter}
    {\normalfont\LARGE\bfseries}
    {\thechapter}
    {1em}
    {}

\titlespacing*{\chapter}
    {0pt}
    {-30pt}
    {18pt}

\titlespacing*{\section}
    {0pt}
    {1.2em}
    {0.5em}

\titlespacing*{\subsection}
    {0pt}
    {1em}
    {0.4em}

% ============================================================
% TEXTE
% ============================================================

\onehalfspacing

\setlength{\parindent}{0pt}
\setlength{\parskip}{0.5em}

% ============================================================
% COULEURS
% ============================================================

\definecolor{ESRFblue}{RGB}{30,80,150}

% Hyperref à la fin
\usepackage{hyperref}

\hypersetup{
    colorlinks=true,
    linkcolor=ESRFblue,
    citecolor=ESRFblue,
    urlcolor=ESRFblue
}

% ============================================================
% DOCUMENT
% ============================================================

\begin{document}


% ============================================================
% PAGE DE GARDE
% ============================================================

\begin{titlepage}

\centering

{\Large European Synchrotron Radiation Facility}\\[0.4cm]
{\large Grenoble, France}

\vfill

{\Huge\bfseries High Performance Algorithms\\
for Tomography}

\vspace{0.7cm}

{\Large Optimisation d'une correction de diffusion optique\\
pour la tomographie X}


\vspace{0.5cm}

{\large Juin -- août 2026}

\vfill

\vspace{1.2cm}

\begin{center}

\begin{tabular}{rl}

\textbf{Stagiaire :} &
{\Large\bfseries Rana ALOUANI}
\\[0.35cm]

\textbf{Formation :} &
Licence 3 Physique
\\[0.25cm]

\textbf{Université :} &
Université Grenoble Alpes
\\[0.8cm]

\textbf{Encadrant :} &
Alessandro Mirone
\\[0.25cm]

\textbf{Organisme :} &
European Synchrotron Radiation Facility (ESRF)

\end{tabular}

\end{center}

\vfill

Grenoble, 2026

\end{titlepage}

% ============================================================
% SYNTHÈSE
% ============================================================

\chapter*{Synthèse}
\addcontentsline{toc}{chapter}{Synthèse}

Ce stage a porté sur l'étude et l'optimisation d'une correction de diffusion
optique utilisée dans le traitement de données de tomographie X à
l'European Synchrotron Radiation Facility.

Dans un système de détection utilisant un scintillateur, les photons X sont
convertis en lumière visible. Une partie de cette lumière se propage avant
d'atteindre la caméra, ce qui redistribue spatialement le signal produit
par l'échantillon.

La correction de cet effet existait déjà dans la chaîne de traitement,
mais son application à pleine résolution sur de grandes projections et sur
un grand nombre d'angles produisait un coût important en calcul et en
stockage.

L'analyse du noyau de diffusion a montré qu'il contient plusieurs échelles
spatiales. Les composantes de grande portée varient lentement dans l'espace
et peuvent être représentées sur une grille plus grossière. Les
contributions plus locales contiennent au contraire davantage
d'information haute fréquence et nécessitent une résolution plus fine.

Cette observation a conduit à séparer la correction en deux composantes :
une partie basse fréquence calculée à résolution spatiale et angulaire
réduite, et une partie haute fréquence appliquée localement.

La basse fréquence a été développée et validée en Python puis stockée sous
forme de données HDF5 réutilisables pendant la reconstruction. La haute
fréquence a été intégrée dans le code C++ de lecture et de prétraitement des
projections, en utilisant FFTW pour réaliser les convolutions locales.

La configuration finale utilise notamment un paramètre de séparation
$\sigma=8$ pixels, un binning spatial de facteur 4 et un espacement
angulaire de facteur 4.

Les mesures finales montrent une réduction de l'empreinte disque d'environ
4,7 To à 80 Go, une réduction du temps de création du digest de 32 à
4 minutes, et un passage du temps de reconstruction de 118 à 93 minutes.

\tableofcontents
\clearpage

% ============================================================
% CHAPITRE 1
% ============================================================

\chapter{Contexte et problématique}

\section{Contexte du stage}

Ce stage de trois mois à l'ESRF a été réalisé volontairement, en dehors
d'une obligation de formation.

L'objectif était de prolonger une formation en physique par une expérience
plus longue sur un problème scientifique réel, à l'interface entre
modélisation physique, traitement numérique des données et développement
logiciel.

Le travail a progressivement évolué d'une compréhension du phénomène
physique vers l'intégration d'une stratégie d'optimisation dans une base
logicielle existante.

\section{Tomographie X}

La tomographie X permet de reconstruire la structure tridimensionnelle
d'un échantillon à partir d'un ensemble de projections enregistrées pour
différents angles.

La qualité du volume reconstruit dépend directement de la qualité des
projections utilisées.

Les imperfections introduites par le système de détection peuvent donc se
propager jusqu'au volume reconstruit.

\section{Problématique}

L'une de ces imperfections provient de la diffusion optique associée au
système de détection.

La correction de cette diffusion était déjà possible dans la chaîne de
traitement.

Le problème étudié pendant le stage était donc le suivant :

\begin{quote}
\centering
\textit{Comment conserver la correction de diffusion optique tout en
réduisant fortement le coût numérique du traitement à pleine résolution ?}
\end{quote}

% ============================================================
% CHAPITRE 2
% ============================================================

\chapter{Origine de la diffusion et formation de l'image}

\section{Du rayonnement X au signal visible}

Dans une expérience de tomographie X, le détecteur ne mesure pas
nécessairement directement les photons X incidents.

Un scintillateur convertit l'énergie déposée par les rayons X en photons
de lumière visible. Cette lumière est ensuite transmise vers le capteur.

Idéalement, l'information produite en une position donnée devrait rester
localisée. Dans un système réel, une partie de la lumière peut cependant
se propager avant d'être enregistrée.

Une source ponctuelle idéale n'est donc pas enregistrée comme un point
strictement ponctuel : son intensité est redistribuée autour de sa
position d'origine.

\section{Modèle par convolution}

Cette redistribution peut être décrite à l'aide d'un noyau de diffusion,
ou \textit{kernel}.

En notant $I_{\mathrm{true}}$ l'image idéale, $K$ le noyau de diffusion et
$I_{\mathrm{measured}}$ l'image enregistrée, on peut écrire de manière
schématique :

\begin{equation}
I_{\mathrm{measured}}
=
I_{\mathrm{true}} * K,
\label{eq:image_convolution}
\end{equation}

où $*$ représente la convolution.

Le noyau décrit donc la manière dont le signal associé à une source
localisée est distribué autour de sa position initiale.

La convolution applique cette même réponse à l'ensemble des pixels de
l'image.


% ============================================================
% CHAPITRE 3
% ============================================================

% ============================================================
% CHAPITRE 3
% ============================================================

\chapter{Un noyau de diffusion multi-échelle}

\section{Profil spatial du noyau}

Le noyau réel n'est pas caractérisé par une seule longueur spatiale.

Plusieurs contributions coexistent : certaines restent très proches du
centre tandis que d'autres s'étendent sur une distance importante.

Le modèle utilisé décrit le noyau comme la superposition d'une contribution
directe et de plusieurs composantes diffuses possédant chacune une échelle
spatiale caractéristique.

La distance au centre du noyau est définie par :

\begin{equation}
r
=
\sqrt{x^2+y^2}.
\label{eq:radial_distance}
\end{equation}

Pour une composante diffuse $i$, le profil radial peut être écrit sous la
forme :

\begin{equation}
K_i(r)
\propto
F_i
\frac{\exp(-r/S_i)}{r},
\qquad r>0.
\label{eq:kernel_component}
\end{equation}

où $S_i$ représente la longueur caractéristique de la composante et $F_i$
son poids relatif.

Le noyau total peut alors être représenté schématiquement par :

\begin{equation}
K(r)
\propto
\delta(r)
+
\sum_i
F_i
\frac{\exp(-r/S_i)}{r}.
\label{eq:kernel_real}
\end{equation}

où :

\begin{itemize}

    \item $r$ représente la distance au centre du noyau ;

    \item $S_i$ est la longueur caractéristique de la composante $i$ ;

    \item $F_i$ est le poids associé à la composante $i$ ;

    \item $\delta(r)$ représente la contribution directe non diffusée.

\end{itemize}

Au voisinage du centre du noyau, lorsque $r \rightarrow 0$, l'expression
analytique des composantes diffuses contient un terme en $1/r$ et ne doit
donc pas être utilisée directement au point central.

Dans l'implémentation, la contribution diffuse est traitée explicitement
au centre afin d'éviter cette singularité numérique. La contribution
directe au point $r=0$ est représentée séparément par le terme
$\delta(r)$.

La divergence apparente de l'expression analytique en $r=0$ ne doit donc
pas être interprétée comme une divergence physique du signal, mais comme
une conséquence de la forme utilisée pour décrire les composantes diffuses
à distance non nulle du centre.


% ============================================================
% FIGURE 1
% ============================================================

\begin{figure}[H]

    \centering

    \includegraphics[width=0.82\textwidth]
    {figures/01_real_diffusion_profile.png}

    \caption{
    Profil vertical du noyau de diffusion réel en échelle logarithmique.
    Le profil met en évidence la coexistence de contributions locales
    autour du centre et de composantes de plus longue portée formant un
    halo spatial étendu.
    }

    \label{fig:real_diffusion_profile}

\end{figure}

La figure~\ref{fig:real_diffusion_profile} montre directement que le
phénomène est multi-échelle.

Les composantes de courte portée produisent des variations très locales,
tandis que les composantes de longue portée donnent naissance au halo
étendu.

Dans le modèle utilisé pour caractériser la diffusion, cette structure
multi-échelle est décrite par plusieurs longueurs caractéristiques.

Les valeurs du paramètre \texttt{scale\_l} correspondent au vecteur :

\begin{equation}
\mathbf{S}
=
\left(
115.0,\,
34.63,\,
7.414,\,
3.472,\,
1.016
\right).
\label{eq:scale_parameters}
\end{equation}

Les fractions associées aux différentes contributions correspondent au
vecteur :

\begin{equation}
\mathbf{F}
=
\left(
0.1859,\,
0.1590,\,
0.1196,\,
0.1375,\,
0.1971
\right).
\label{eq:fraction_parameters}
\end{equation}

Chaque couple $(S_i,F_i)$ correspond ainsi à une composante du noyau
caractérisée par une portée spatiale et un poids différents.

Les grandes valeurs de $S_i$ décrivent les contributions les plus
étendues, tandis que les petites valeurs correspondent à des contributions
plus localisées autour du centre.

Les valeurs de \texttt{scale\_l} ne doivent pas être interprétées comme des
largeurs gaussiennes $\sigma$. Elles représentent les longueurs
caractéristiques des différentes composantes du modèle de diffusion.

La coexistence de ces différentes échelles spatiales constitue le point
de départ de la stratégie d'optimisation développée dans la suite de ce
travail.

% ============================================================
% CHAPITRE 4
% ============================================================

\chapter{Pourquoi la correction complète est coûteuse}

\section{Déconvolution de référence}

Pour corriger la diffusion, la méthode de référence réalise une
déconvolution.

Dans le domaine de Fourier, on peut écrire de manière schématique :

\begin{equation}
\widehat{I}_{\mathrm{corr}}
=
\frac{\widehat{I}_{\mathrm{meas}}}
{\widehat{K}},
\label{eq:deconv}
\end{equation}

où :

\begin{itemize}

    \item $\widehat{I}_{\mathrm{corr}}$ est la transformée de Fourier de la
    projection corrigée ;

    \item $\widehat{I}_{\mathrm{meas}}$ est la transformée de la projection
    mesurée ;

    \item $\widehat{K}$ est la transformée de Fourier du noyau de diffusion.

\end{itemize}

En pratique, cette correction nécessite plusieurs opérations successives.
Chaque projection est d'abord étendue par padding afin de limiter les
effets de bord. Une transformée de Fourier est ensuite calculée afin de
passer dans le domaine fréquentiel, où la correction est appliquée à partir
du noyau de diffusion. Une transformée de Fourier inverse permet ensuite de
revenir dans l'espace réel, avant de recadrer l'image pour retrouver la
zone utile de la projection.

Cette procédure devient coûteuse car elle doit être répétée pour un très
grand nombre de projections, chacune étant traitée à pleine résolution.
Chaque projection nécessite donc des opérations de padding, de transformée
de Fourier, de correction dans le domaine fréquentiel puis de transformée
inverse sur des tableaux de grande taille. À l'échelle d'une acquisition
tomographique complète, cela représente un volume important de données à
manipuler et un coût de calcul élevé.

Une question se pose alors : si une partie de la correction varie seulement
lentement dans l'espace et entre des projections voisines, est-il réellement
nécessaire de la calculer partout et à pleine résolution ?

\vspace{0.8em}

L'idée consiste donc à ne plus traiter l'ensemble de la correction de la
même manière, mais à introduire une
\textbf{\textcolor{ESRFblue}{séparation basse fréquence / haute fréquence}} :
une composante basse fréquence, lente et étendue, et une composante haute
fréquence, plus locale et détaillée.


% ============================================================
% CHAPITRE 5
% ============================================================

\chapter{Séparation basse fréquence et haute fréquence}

\section{Définition de la correction}

On définit tout d'abord la correction $Q$ comme la différence entre le
signal mesuré et le signal entièrement corrigé :

\begin{equation}
Q
=
I_{\mathrm{meas}}
-
I_{\mathrm{corr}},
\end{equation}

où $I_{\mathrm{meas}}$ désigne la projection mesurée par le détecteur et
$I_{\mathrm{corr}}$ la projection obtenue après application de la correction
complète de référence.

L'idée consiste ensuite à ne plus traiter cette correction comme un seul
objet, mais à la séparer en deux contributions complémentaires : une
composante basse fréquence et une composante haute fréquence.

La composante basse fréquence est obtenue en appliquant un filtre gaussien
à la correction :

\begin{equation}
Q_{\mathrm{LF}}
=
G_\sigma * Q,
\label{eq:lf}
\end{equation}

où $G_\sigma$ désigne un filtre gaussien de largeur $\sigma$.

Dans le travail présenté ici, la valeur

\begin{equation}
\sigma = 8\ \mathrm{pixels}
\end{equation}

a été retenue. Cette largeur permet de conserver dans $Q_{\mathrm{LF}}$ les
variations lentes que l'on souhaite traiter comme une composante globale,
tout en séparant suffisamment les structures plus locales. Le choix de
$\sigma$ fixe donc la frontière entre les deux échelles de correction.

La composante haute fréquence n'est pas calculée indépendamment. Elle est
simplement définie comme ce qui reste après soustraction de la basse
fréquence :

\begin{equation}
Q_{\mathrm{HF}}
=
Q
-
Q_{\mathrm{LF}}.
\label{eq:hf}
\end{equation}

Par construction, les deux contributions redonnent donc exactement la
correction initiale :

\begin{equation}
Q
=
Q_{\mathrm{LF}}
+
Q_{\mathrm{HF}}.
\label{eq:sum_lf_hf}
\end{equation}

Cette propriété est importante : la séparation ne modifie pas à ce stade
la correction de référence. Elle constitue uniquement une décomposition
de celle-ci en deux contributions présentant des propriétés spatiales
différentes.

La composante basse fréquence contient principalement les variations
lentes et spatialement étendues, tandis que la composante haute fréquence
conserve les structures plus localisées.


La même construction peut être appliquée au \textbf{noyau de la correction}.
C'est cette décomposition du noyau qui est représentée sur la
figure~\ref{fig:low_high_separation}. La figure ne montre donc pas directement
la décomposition d'une projection corrigée, mais la séparation du noyau de
correction en une contribution basse fréquence et une contribution haute
fréquence.



% ============================================================
% FIGURE 2
% ============================================================

% FIGURE DE LA SLIDE 5
%
% Fichier de la présentation :
% kernel_correction_low_high_sigma8_ZOOM.png
%
% Nom conseillé dans Overleaf :
% figures/02_low_high_separation.png

\begin{figure}[H]
    \centering
    \includegraphics[width=0.9\textwidth]
{figures/02_correction_kernel_profile.png}
    \caption{
    Décomposition du noyau de correction en une composante basse fréquence
    et une composante haute fréquence pour $\sigma=8$ pixels. La composante
    basse fréquence conserve la structure spatialement étendue du noyau,
    tandis que la composante haute fréquence est davantage concentrée
    autour de son centre. Les échelles d'affichage sont adaptées
    indépendamment afin de rendre visible la composante haute fréquence.
    }
    \label{fig:low_high_separation}
\end{figure}


Cette décomposition fait apparaître deux propriétés différentes qui peuvent
être exploitées numériquement : la basse fréquence est lisse et peut
potentiellement être représentée à plus faible résolution, tandis que la
haute fréquence est plus localisée et peut potentiellement être traitée sur
une région limitée. La stratégie d'optimisation consiste donc à adapter le
coût du calcul aux propriétés propres de chacune de ces deux composantes.


% ============================================================
% CHAPITRE 6
% ============================================================

\chapter{Réduction du coût de la basse fréquence}

\section{Réduction spatiale}

La première optimisation concerne la composante basse fréquence.
Comme elle varie lentement dans l'espace, il n'est pas nécessaire de la
représenter avec la même densité de pixels que la correction complète.

On peut donc commencer par réduire sa résolution spatiale avant de la traiter. 


Pour une projection de dimensions $N_x\times N_y$, un binning spatial
de facteur $b$ dans les deux directions conduit à une représentation de
dimensions :

\begin{equation}
\frac{N_x}{b}
\times
\frac{N_y}{b}.
\end{equation}

Le nombre total d'échantillons spatiaux devient alors :

\begin{equation}
\frac{N_xN_y}{b^2}.
\end{equation}

Avec $b=4$, cela correspond à une division par 16 du nombre
d'échantillons spatiaux.



Ce choix reste compatible avec la séparation fréquentielle utilisée ici.
Pour $\sigma=8$ pixels, les structures conservées dans la composante basse
fréquence sont déjà lissées sur plusieurs pixels. Après un binning spatial
de facteur 4, cette échelle caractéristique correspond encore à environ
deux pixels dans l'image réduite. La taille des données peut donc être
fortement diminuée sans supprimer les variations lentes que l'on cherche
précisément à conserver.

Cette réduction d'un facteur 16 concerne cependant uniquement le nombre
d'échantillons spatiaux. Elle ne signifie pas que l'algorithme complet est
automatiquement 16 fois plus rapide. Le gain réel dépend de l'ensemble de
l'implémentation, notamment du coût des FFT, des accès mémoire et des
entrées-sorties.

\section{Réduction angulaire}


La réduction spatiale n'est toutefois pas la seule possibilité. La
composante basse fréquence évolue également relativement lentement d'une
projection à la suivante. Cette propriété suggère une seconde réduction,
cette fois dans la direction angulaire : au lieu de calculer directement
la basse fréquence pour chaque projection, il devient possible de ne la
calculer que pour certaines projections puis d'interpoler les valeurs
intermédiaires.

Pour vérifier cette idée, un test a été réalisé sur la projection 102.
La correction basse fréquence de cette projection n'est pas utilisée pour
construire la prédiction. Elle est conservée uniquement comme référence.
Les corrections basse fréquence sont calculées pour les projections 100
et 104, puis réduites spatialement d'un facteur 4. La valeur correspondant
à la projection 102 est ensuite obtenue par interpolation entre ces deux
projections voisines.

Dans ce cas, le poids d'interpolation vaut :

\begin{equation}
t
=
\frac{102-100}{104-100}
=
\frac{1}{2}.
\end{equation}

L'interpolation tient également compte d'une normalisation par le courant
du faisceau, car l'intensité globale peut varier d'une projection à
l'autre.

Afin de comparer cette approximation à la référence, la
figure~\ref{fig:lf_interpolation} présente à gauche la véritable correction
basse fréquence de la projection 102, au centre la correction prédite à
partir des projections 100 et 104, et à droite leur différence.



% ============================================================
% FIGURE 3
% ============================================================

% FIGURE DE LA SLIDE 7
%
% Elle doit contenir :
% - vraie LF projection 102
% - LF interpolée depuis 100 et 104
% - résidu
%
% Nom :
% figures/03_lf_interpolation.png

\begin{figure}[H]
    \centering
    \includegraphics[width=0.95\textwidth] {figures/03_lf_spatial_angular_interpolation.png}
    \caption{
    Validation de la réduction spatiale et angulaire de la composante basse
    fréquence pour la projection 102. À gauche : correction basse fréquence
    de référence. Au centre : prédiction obtenue à partir des projections
    100 et 104 avec un binning spatial de facteur 4 et une interpolation
    angulaire. À droite : résidu entre la référence et la prédiction.
    La carte du résidu utilise une échelle divergente centrée autour de zéro.
    }
    \label{fig:lf_interpolation}
\end{figure}


Les structures de grande échelle sont visuellement très proches entre la
référence et la prédiction. Le résidu peut sembler marqué en raison de
l'échelle de couleurs utilisée, mais son amplitude reste faible par rapport
à celle de la correction basse fréquence elle-même.

Les écarts-types mesurés sont :

\begin{align}
\mathrm{std}(Q_{\mathrm{LF,true}}) &= 137.6,\\
\mathrm{std}(Q_{\mathrm{LF,interp}}) &= 136.9,\\
\mathrm{std}(\mathrm{residual}) &= 3.81.
\end{align}

L'écart-type du résidu représente donc environ $2.8\,\%$ de celui de la
correction basse fréquence de référence.

Ce test montre que la composante lisse peut être approximée en combinant
une réduction spatiale et une réduction angulaire, tout en conservant les
structures basse fréquence que l'on souhaite préserver. Cette validation
justifie ensuite l'utilisation de ces deux réductions dans la configuration
finale du pipeline.

% ============================================================
% CHAPITRE 7
% ============================================================

\chapter{Traitement local de la haute fréquence}

\section{Exploitation de la localité spatiale}

Contrairement à la composante basse fréquence, la composante haute
fréquence ne peut pas être fortement réduite spatialement. Elle contient
précisément les structures fines que l'on souhaite préserver, et un binning
trop important supprimerait une partie de cette information.

L'optimisation doit donc exploiter une autre propriété de cette
composante : sa \textbf{\textcolor{ESRFblue}{localité spatiale}}.

La figure~\ref{fig:hf_profile} représente la magnitude normalisée du noyau
de correction haute fréquence en fonction de la distance verticale à son
centre. L'axe vertical est représenté en échelle logarithmique afin de
rendre visibles des amplitudes qui s'étendent sur plusieurs ordres de
grandeur.

Au voisinage du centre, l'amplitude de la composante haute fréquence est
importante. En revanche, lorsqu'on s'éloigne du centre du noyau, cette
amplitude décroît très rapidement. Le graphe montre donc que l'essentiel de
la contribution haute fréquence est concentré dans une région spatiale
relativement limitée.

Cette observation est particulièrement intéressante du point de vue
numérique. Si la contribution haute fréquence devient négligeable loin du
centre, il n'est pas nécessaire de la calculer sur toute la hauteur de la
projection. Il suffit de conserver une marge autour de la région qui doit
effectivement être corrigée.


% FIGURE 4
% ============================================================

% FIGURE DE LA SLIDE 8
%
% Fichier original :
% HF_kernel_vertical_profile_sigma8_FINAL.png
%
% Nom :
% figures/04_hf_local_profile.png

\begin{figure}[H]
    \centering
    \includegraphics[width=0.90\textwidth]
    {figures/04_hf_local_profile.png}
    \caption{
    Magnitude normalisée du noyau de correction haute fréquence en fonction
    de la distance verticale à son centre, représentée en échelle
    logarithmique. L'amplitude décroît de plusieurs ordres de grandeur
    lorsque l'on s'éloigne du centre. Les différentes marges testées
    permettent d'évaluer la taille de voisinage nécessaire pour reproduire
    correctement la contribution haute fréquence.
    }
    \label{fig:hf_profile}
\end{figure}


Afin de déterminer jusqu'à quelle distance cette contribution doit être
conservée, plusieurs marges locales ont été testées autour de la région
utile.

Les erreurs obtenues pour la composante haute fréquence sont résumées dans
le tableau suivant :

\begin{center}
\begin{tabular}{cc}
\toprule
\textbf{Marge locale} & \textbf{RMSE haute fréquence}\\
\midrule
16 px  & 0.828\\
32 px  & 0.322\\
64 px  & 0.150\\
128 px & 0.0225\\
\bottomrule
\end{tabular}
\end{center}

Avec une marge de 16 pixels, l'erreur reste encore relativement importante,
avec une RMSE d'environ 0.83. Lorsque la marge est portée à 32 pixels,
l'erreur diminue à environ 0.32. À 64 pixels, elle atteint environ 0.15,
puis seulement 0.0225 avec une marge de 128 pixels.

Cette diminution rapide confirme quantitativement ce que suggère le profil
du noyau : la majeure partie de la contribution haute fréquence est
concentrée près du centre.

Il n'est donc pas nécessaire de traiter toute la hauteur de la projection
pour retrouver correctement la haute fréquence. Un voisinage local autour
de la région utile suffit à en capturer l'essentiel.


% ============================================================
% CHAPITRE 8
% ============================================================
\chapter{Pipeline multirésolution complet}

\section{Combinaison des deux approximations}

À ce stade, les deux idées d'optimisation ont été validées séparément.
D'une part, la composante basse fréquence peut être approchée à résolution
spatiale réduite et interpolée angulairement. D'autre part, la composante
haute fréquence peut être calculée localement, car l'essentiel de sa
contribution est concentré dans un voisinage limité autour de la région
utile.

L'étape suivante consiste donc à combiner ces deux approximations pour
reconstruire une approximation de la correction complète.

On définit ainsi :

\begin{equation}
Q_{\mathrm{approx}}
=
Q_{\mathrm{LF,pred}}
+
Q_{\mathrm{HF,local}},
\label{eq:qapprox}
\end{equation}

où :

\begin{itemize}

    \item $Q_{\mathrm{LF,pred}}$ est la composante basse fréquence prédite
    à partir d'un calcul à résolution réduite, avec binning spatial et
    interpolation angulaire ;

    \item $Q_{\mathrm{HF,local}}$ est la composante haute fréquence calculée
    localement autour de la région utile.

\end{itemize}

La projection corrigée approchée s'écrit alors :

\begin{equation}
I_{\mathrm{corr,approx}}
=
I_{\mathrm{meas}}
-
Q_{\mathrm{approx}}.
\label{eq:icorrapprox}
\end{equation}

Dans la configuration étudiée ici, la composante basse fréquence est
obtenue avec un binning spatial de facteur 4 et un binning angulaire de
facteur 4, tandis que la composante haute fréquence est calculée avec une
marge locale de 64 pixels.

La figure~\ref{fig:multiscale_pipeline} montre comment ces deux
approximations sont combinées sur le cas de la projection 102. En haut à
gauche apparaît la véritable correction basse fréquence de référence pour
cette projection, tandis qu'en haut à droite figure sa prédiction obtenue
à partir des projections voisines avec réduction spatiale et angulaire.
En bas à gauche est représentée la correction haute fréquence calculée
localement avec une marge de 64 pixels. L'addition de ces deux
contributions fournit alors l'approximation complète de la correction.
Enfin, le panneau en bas à droite montre le résidu final par rapport à la
correction de référence.


% ============================================================
% FIGURE 5
% ============================================================

% FIGURE DE LA SLIDE 9
%
% Fichier original :
% FINAL_pipeline_multiscale_margin64.png
%
% Nom :
% figures/05_multiscale_pipeline.png

\begin{figure}[H]
    \centering
    \includegraphics[width=0.95\textwidth]
    {figures/05_multiscale_pipeline.png}
    \caption{
    Combinaison des deux approximations dans le pipeline multirésolution.
    En haut à gauche : correction basse fréquence de référence pour la
    projection 102. En haut à droite : prédiction basse fréquence obtenue
    avec un binning spatial de facteur 4 et un binning angulaire de facteur
    4. En bas à gauche : correction haute fréquence calculée localement avec
    une marge de 64 pixels. En bas à droite : résidu final entre
    l'approximation complète et la correction de référence.
    }
    \label{fig:multiscale_pipeline}
\end{figure}


Le point important est que le résidu final apparaît beaucoup plus lisse que
la composante haute fréquence elle-même. Cela signifie que l'erreur liée à
la localisation de la haute fréquence a déjà été fortement réduite.

Cette observation est cohérente avec les résultats quantitatifs obtenus
précédemment. Pour une marge de 64 pixels, l'erreur propre à la
localisation de la haute fréquence est caractérisée par une RMSE
d'environ 0.15, tandis que l'erreur finale de reconstruction de la
correction complète est d'environ 2.23.

Autrement dit, à ce stade, la limitation principale ne provient plus de la
troncature locale de la haute fréquence. L'erreur restante est dominée par
l'approximation de la composante basse fréquence, c'est-à-dire par la
réduction spatiale et l'interpolation angulaire utilisées pour prédire
$Q_{\mathrm{LF,pred}}$.

Cette étape montre donc que les deux contributions peuvent être recombinées
de manière cohérente, et que la stratégie multirésolution permet de
reproduire la correction complète avec une erreur principalement contrôlée
par la branche basse fréquence.


% ============================================================
% CHAPITRE 9
% ============================================================

\chapter{Analyse de l'erreur restante}

Les résultats quantitatifs permettent d'identifier plus précisément quelle
composante limite désormais la précision de la méthode.

Lorsque la marge utilisée pour le calcul de la composante haute fréquence
augmente, l'approximation de cette branche s'améliore fortement. Une marge
plus grande permet en effet de conserver une portion plus importante du
noyau haute fréquence et donc de réduire l'erreur associée à sa
localisation.

Les valeurs obtenues sont résumées dans le tableau suivant :

\begin{center}
\begin{tabular}{ccc}
\toprule
\textbf{Marge HF} &
\textbf{RMSE HF} &
\textbf{RMSE finale}\\
\midrule
16 px  & 0.828  & 2.349\\
32 px  & 0.322  & 2.243\\
64 px  & 0.150  & 2.231\\
128 px & 0.0225 & 2.226\\
192 px & 0      & 2.226\\
\bottomrule
\end{tabular}
\end{center}

La RMSE associée à la haute fréquence diminue donc très rapidement lorsque
la marge augmente : elle passe d'environ 0.83 pour une marge de 16 pixels
à 0.15 pour 64 pixels, puis devient presque nulle lorsque l'on utilise une
région suffisamment grande.

En revanche, l'erreur finale évolue beaucoup moins. Elle passe de 2.349
pour une marge de 16 pixels à environ 2.23 dès que la marge atteint
64 pixels, puis reste pratiquement constante lorsque la marge est encore
augmentée.

Au-delà d'environ 64 à 128 pixels, augmenter davantage la marge haute
fréquence n'apporte donc presque plus d'amélioration au résultat final.

% ============================================================
% FIGURE 6
% ============================================================

\begin{figure}[H]
    \centering
    \includegraphics[width=0.90\textwidth]
    {figures/06_hf_margin_error.png}
    \caption{
    Évolution de l'erreur en fonction de la marge utilisée pour le calcul
    local de la composante haute fréquence. L'erreur haute fréquence diminue
    fortement lorsque la marge augmente, tandis que l'erreur finale se
    stabilise autour de 2.23 à partir d'une marge d'environ 64 pixels.
    }
    \label{fig:hf_margin_error}
\end{figure}

Ce comportement permet d'identifier l'origine de l'erreur restante.
Lorsque la marge haute fréquence devient suffisamment grande, l'erreur liée
à la localisation de cette composante devient faible devant l'erreur
totale.

La RMSE finale, proche de 2.23, est alors très proche de l'erreur associée
à la prédiction basse fréquence, d'environ 2.226 dans ce test. Le facteur
limitant ne se situe donc plus dans la branche haute fréquence, mais dans
l'approximation de la composante basse fréquence.

Autrement dit, la localisation de la haute fréquence fonctionne
suffisamment bien. Une amélioration supplémentaire de la précision du
pipeline devra désormais porter principalement sur la prédiction basse
fréquence, notamment sur la réduction spatiale et l'interpolation
angulaire.


% ============================================================
% CHAPITRE 10
% ============================================================

\chapter{Implémentation dans la chaîne de reconstruction}

\section{Nouvelle organisation de la correction}

La séparation entre basse et haute fréquence ne modifie pas seulement la
manière de calculer la correction. Elle permet également de modifier
l'endroit où chaque contribution est appliquée dans la chaîne de
reconstruction.

Dans l'approche précédente, la correction complète devait être réalisée
avant la création du digest. Sa composante de longue portée nécessite en
effet des informations provenant d'une région importante de la
radiographie, ce qui impose de travailler sur des données beaucoup plus
étendues que la seule zone finalement reconstruite.

Après séparation des échelles, cette contrainte ne concerne plus toute la
correction. La composante basse fréquence reste globale et doit donc être
préparée en amont. À l'inverse, la composante haute fréquence est
suffisamment localisée pour être appliquée plus tard, directement pendant
le traitement du volume, sur la région déjà disponible dans la marge de
lecture.

La nouvelle organisation est donc la suivante :

\begin{center}
\textbf{Basse fréquence : préparation en amont à résolution réduite}
\\[0.3em]
+
\\[0.3em]
\textbf{Haute fréquence : application locale pendant la lecture}
\end{center}

Cette organisation permet d'éviter de produire et de stocker une correction
complète à pleine résolution pour toutes les projections.

\section{Pourquoi utiliser Python pour la basse fréquence et C++ pour la haute fréquence ?}

Les deux composantes n'interviennent pas au même moment et n'ont pas les
mêmes contraintes numériques. Il n'était donc pas nécessaire de les
implémenter dans le même environnement.

La composante basse fréquence correspond à une étape de préparation des
données. Elle doit construire le modèle de diffusion, manipuler les
paramètres expérimentaux, effectuer les opérations de réduction spatiale et
angulaire, puis produire des données intermédiaires qui seront réutilisées
pendant la reconstruction.

Python est bien adapté à cette phase grâce aux outils de calcul scientifique
déjà utilisés pendant le développement, notamment NumPy, SciPy, h5py et
pyFFTW. Il permet également de conserver une implémentation proche des
scripts utilisés pendant la phase de prototypage.

La composante haute fréquence répond à une contrainte différente. Elle doit
être appliquée au moment où les projections sont déjà lues par
\texttt{nabuxx} et présentes en mémoire. L'intégrer directement dans le code
C++ évite donc de créer une nouvelle étape intermédiaire ou de réécrire les
projections sur disque avant la reconstruction.

Le choix Python/C++ ne correspond ainsi pas simplement à une préférence de
langage : il suit la position naturelle de chaque composante dans le
pipeline.

\section{Préparation de la basse fréquence en Python}

\subsection{Entrées du calcul}

Le programme Python principal utilisé pour cette étape est :

\begin{center}
\texttt{diffusion\_correction.py}
\end{center}

Il utilise les projections expérimentales, le dark et les paramètres du
modèle de diffusion. Les principaux fichiers d'entrée sont notamment :

\begin{itemize}
    \item \texttt{projections.h5} ;
    \item \texttt{dark.h5} ;
    \item \texttt{diffusion\_parameters.json}.
\end{itemize}

Le programme récupère les dimensions des projections ainsi que le courant
associé à chaque image. Il applique ensuite les facteurs de réduction
spatiale et angulaire définis précédemment.

Par exemple, les dimensions réduites sont calculées dans le programme à
partir des facteurs \texttt{bin\_proj} et \texttt{bin\_x} :

\begin{verbatim}
bin_proj = args.diffusion_bin_proj
bin_x = args.diffusion_bin_x

nprojs_binned = (nprojs + bin_proj - 1) // bin_proj
dim_z_binned = (dim_z + bin_x - 1) // bin_x
dim_x_binned = (dim_x + bin_x - 1) // bin_x
\end{verbatim}

Cette étape traduit directement dans le code la réduction spatiale et
angulaire étudiée précédemment.

\subsection{Calcul et stockage de la correction réduite}

La correction n'est donc calculée que pour les projections retenues par le
binning angulaire et sur une grille spatiale réduite.

La composante basse fréquence obtenue doit ensuite rester disponible pour
la reconstruction. Elle est enregistrée dans le fichier :

\begin{center}
\texttt{diffusion\_correction.h5}
\end{center}

Le fichier contient notamment trois datasets :

\begin{itemize}
    \item \texttt{difference}, qui contient les valeurs de correction ;
    \item \texttt{binnings}, qui conserve les facteurs de réduction
    utilisés ;
    \item \texttt{current}, qui conserve le courant associé aux projections.
\end{itemize}

La création de ces datasets apparaît directement dans le code :

\begin{verbatim}
f_out.create_dataset(
    "difference",
    expected_shape,
    dtype="f4"
)

f_out.create_dataset(
    "binnings",
    data=np.asarray([bin_proj, bin_x], dtype="i4")
)

f_out.create_dataset(
    "current",
    shape=(nprojs_binned,),
    dtype="f4"
)
\end{verbatim}

Le fichier HDF5 constitue ainsi l'interface entre la phase de préparation
Python et la reconstruction réalisée ensuite par \texttt{nabuxx}.

\subsection{Réutilisation pendant la reconstruction}

Lorsqu'un bloc de projections doit être traité, il n'est pas nécessaire de
charger toute la correction basse fréquence en mémoire.

\texttt{nabuxx} récupère uniquement la portion du volume de correction
nécessaire au bloc actuellement traité. Les valeurs enregistrées à
résolution réduite sont ensuite interpolées suivant les dimensions
nécessaires pour retrouver la résolution des projections utilisées par la
reconstruction.

L'interpolation intervient donc suivant trois directions :

\begin{itemize}
    \item l'indice de projection ;
    \item la direction verticale ;
    \item la direction horizontale.
\end{itemize}

Cette étape relie la représentation économique construite en Python aux
données brutes actuellement présentes dans le pipeline.

\section{Application locale de la haute fréquence en C++}

\subsection{Préparation du kernel haute fréquence}

La haute fréquence doit être appliquée dans le code C++, mais il n'est pas
utile de reconstruire le modèle physique du kernel à chaque lecture d'une
projection.

Le kernel haute fréquence est donc préparé une seule fois en Python puis
enregistré dans :

\begin{center}
\texttt{kernel.h5}
\end{center}

Ce fichier contient un dataset bidimensionnel nommé :

\begin{center}
\texttt{kernel}
\end{center}

Le C++ peut ainsi charger directement le kernel déjà calculé et se
concentrer uniquement sur son application aux données.

\subsection{Chargement dans \texttt{nabuxx}}

Le kernel est chargé au démarrage de la correction à partir du fichier
HDF5. Le code utilise HighFive pour accéder au dataset :

\begin{verbatim}
h5::File kernel_file(
    (fs::path(mypath) / "kernel.h5").string(),
    h5::File::ReadOnly
);

h5::DataSet kernel_dataset =
    kernel_file.getDataSet("kernel");
\end{verbatim}

Ses dimensions sont vérifiées avant d'être copiées dans la structure
utilisée par le calcul C++.


\subsection{Fonction dédiée à la correction haute fréquence}

La correction locale a été isolée dans une fonction dédiée dans
\texttt{etfscan.cpp} :

\begin{verbatim}
void apply_high_frequency_convolution(
    xt::xtensor<float, 3>& raw_data,
    const xt::xarray<float>& kernel
)
\end{verbatim}

Cette fonction reçoit deux éléments :

\begin{itemize}
    \item le bloc de projections actuellement présent dans
    \texttt{raw\_data} ;
    \item le kernel haute fréquence pré-calculé.
\end{itemize}

Son rôle est volontairement limité : appliquer la contribution haute
fréquence aux données déjà chargées en mémoire puis réinjecter le résultat
dans le même bloc.

Cette séparation permet de conserver une fonction identifiable et testable,
sans mélanger directement le calcul de convolution avec toute la logique de
lecture des fichiers HDF5.


\subsection{Convolution locale avec FFTW}

La convolution haute fréquence est réalisée à l'aide de FFTW.

Une convolution bidimensionnelle directe entre une projection et un kernel
peut nécessiter un grand nombre d'opérations. Le calcul est donc réalisé
dans le domaine de Fourier en utilisant le théorème de convolution :

\begin{equation}
I*K
=
\mathcal{F}^{-1}
\left[
\mathcal{F}(I)\,
\mathcal{F}(K)
\right],
\end{equation}

où $I$ représente la région de projection traitée et $K$ le kernel haute
fréquence.

Le travail réalisé ne consistait pas à réimplémenter l'algorithme FFT
lui-même. FFTW fournit les transformées nécessaires. L'intégration porte
principalement sur la préparation des données, les dimensions de padding,
la gestion des buffers, la multiplication des spectres, la transformée
inverse et le recadrage du résultat.

Pour une image de dimensions $H_I\times W_I$ et un kernel de dimensions
$H_K\times W_K$, les dimensions nécessaires à la convolution linéaire sont :

\begin{align}
H_P &= H_I + H_K - 1,\\
W_P &= W_I + W_K - 1.
\end{align}

Dans le code :

\begin{verbatim}
const std::size_t padded_height =
    image_height + kernel_height - 1;

const std::size_t padded_width =
    image_width + kernel_width - 1;
\end{verbatim}

Après la transformée inverse, le résultat est normalisé puis recadré afin
de retrouver les dimensions attendues par la suite du pipeline.

\section{Flux final des données}

L'implémentation finale peut donc être résumée de la manière suivante.

Les fichiers \texttt{projections.h5}, \texttt{dark.h5} et
\texttt{diffusion\_parameters.json} sont d'abord utilisés par le programme
Python pour préparer la composante basse fréquence à résolution spatiale et
angulaire réduite.

Cette correction est sauvegardée dans :

\begin{center}
\texttt{diffusion\_correction.h5}.
\end{center}

En parallèle, le kernel haute fréquence est pré-calculé en Python et
enregistré dans :

\begin{center}
\texttt{kernel.h5}.
\end{center}

Pendant la reconstruction, \texttt{nabuxx} charge uniquement la partie
nécessaire de la basse fréquence, l'interpole pour retrouver les valeurs
correspondant au bloc de projections courant, puis applique localement la
haute fréquence directement en C++.

On obtient ainsi le flux suivant :

\begin{figure}[H]
\centering

\resizebox{0.88\textwidth}{!}{
\begin{tikzpicture}[
    >=Latex,
    node distance=0.8cm and 1.2cm,
    box/.style={
        draw=ESRFblue,
        rounded corners=3pt,
        thick,
        align=center,
        minimum height=0.9cm,
        text width=4.5cm,
        fill=blue!3,
        inner sep=5pt
    },
    bigbox/.style={
        draw=ESRFblue,
        rounded corners=3pt,
        thick,
        align=center,
        minimum height=0.9cm,
        text width=5.8cm,
        fill=blue!3,
        inner sep=5pt
    },
    filebox/.style={
        draw=black!45,
        rounded corners=3pt,
        align=center,
        minimum height=0.8cm,
        text width=4.8cm,
        fill=black!2,
        inner sep=4pt
    },
    arrow/.style={
        ->,
        thick
    }
]

% Entrées
\node[bigbox] (input)
{\texttt{projections.h5} + \texttt{dark.h5}\\
+ \texttt{diffusion\_parameters.json}};

% Branches
\node[box, below left=1.3cm and 1.8cm of input] (lf)
{\textbf{Python}\\
Basse fréquence réduite};

\node[box, below right=1.3cm and 1.8cm of input] (hf)
{\textbf{Python}\\
Préparation du kernel HF};

% Fichiers
\node[filebox, below=0.75cm of lf] (lfout)
{\texttt{diffusion\_correction.h5}};

\node[filebox, below=0.75cm of hf] (hfout)
{\texttt{kernel.h5}};

% Centre
\node[bigbox, below=3.6cm of input] (cpp)
{\textbf{C++ / \texttt{nabuxx} / \texttt{etfscan.cpp}}};

\node[bigbox, below=0.8cm of cpp] (process)
{Interpolation LF\\
+ correction HF locale};

\node[filebox, below=0.8cm of process] (corr)
{Projections corrigées};

\node[filebox, below=0.8cm of corr] (recons)
{Reconstruction tomographique};

% Flèches propres
\draw[arrow] (input.south west) |- (lf.north);
\draw[arrow] (input.south east) |- (hf.north);

\draw[arrow] (lf.south) -- (lfout.north);
\draw[arrow] (hf.south) -- (hfout.north);

\draw[arrow] (lfout.south east) |- (cpp.west);
\draw[arrow] (hfout.south west) |- (cpp.east);

\draw[arrow] (cpp.south) -- (process.north);
\draw[arrow] (process.south) -- (corr.north);
\draw[arrow] (corr.south) -- (recons.north);

\end{tikzpicture}
}

\caption{
Organisation du pipeline après séparation basse fréquence / haute fréquence.
La basse fréquence est préparée en Python à résolution réduite et stockée
dans \texttt{diffusion\_correction.h5}. Le kernel haute fréquence est
pré-calculé séparément et stocké dans \texttt{kernel.h5}. Les deux éléments
sont ensuite utilisés dans \texttt{nabuxx}, où la basse fréquence est
interpolée et la correction haute fréquence appliquée localement avant la
reconstruction.
}

\label{fig:pipeline_lf_hf}

\end{figure}

Cette architecture associe donc chaque composante au traitement numérique
le mieux adapté à ses propriétés : une préparation réduite et réutilisable
pour la basse fréquence, et une application locale directement sur les
données en mémoire pour la haute fréquence.





% ============================================================
% CHAPITRE 11
% ============================================================

\chapter{Validation sur reconstruction expérimentale}

Après avoir validé séparément les différentes composantes de la correction,
il reste à vérifier leur effet sur le résultat final de la reconstruction
tomographique.

Toutes les images présentées dans cette section correspondent à la
\textbf{même coupe reconstruite}, avec le \textbf{même recadrage} et la
\textbf{même plage de niveaux de gris}. Elles peuvent donc être comparées
directement sans que les différences observées proviennent d'un changement
d'affichage.

\section{Effet progressif de la correction}

La première comparaison permet de visualiser progressivement l'effet des
différentes étapes de correction.

% ============================================================
% FIGURE 7
% ============================================================

% Image à construire avec trois panneaux :
% (a) aucune déconvolution
% (b) correction basse fréquence seule
% (c) nouvelle correction LF + HF
%
% Même coupe, même crop, même grayscale.
%
% Nom conseillé :
% figures/07_reconstruction_progression.png

\begin{figure}[H]
\centering

\begin{tikzpicture}[
    >=Latex,
    node distance=0.8cm,
    arrow/.style={->, thick, draw=ESRFblue}
]

% Images
\node (img1) {
    \includegraphics[width=0.28\textwidth]
    {figures/no deconv.png}
};

\node[right=1.4cm of img1] (img2) {
    \includegraphics[width=0.28\textwidth]
    {figures/only lf.png}
};

\node[right=1.4cm of img2] (img3) {
    \includegraphics[width=0.28\textwidth]
    {figures/correc new.png}
};

% Flèches
\draw[arrow] (img1.east) --
node[midway, above=4pt, font=\scriptsize, align=center]
{LF\\correction}
(img2.west);

\draw[arrow] (img2.east) --
node[midway, above=4pt, font=\scriptsize, align=center]
{+ HF local\\correction}
(img3.west);

% Labels sous les images
\node[below=0.2cm of img1]
{\small No deconvolution};

\node[below=0.2cm of img2]
{\small Low-frequency correction only};

\node[below=0.2cm of img3]
{\small Final LF + HF correction};

\end{tikzpicture}

\caption{
Effet progressif de la correction sur une même coupe reconstruite.
De gauche à droite : sans déconvolution, avec correction basse fréquence
seule, puis avec la correction finale LF + HF. Les trois images utilisent
le même recadrage et la même plage de niveaux de gris.
}

\label{fig:reconstruction_progression}

\end{figure}



Dans la première image, aucune déconvolution n'est appliquée. Elle
correspond donc à la reconstruction obtenue avant correction de la
diffusion optique.

Dans la deuxième image, seule la composante basse fréquence est appliquée.
Une partie de l'effet de la correction est déjà retrouvée, mais cette
composante contient volontairement uniquement la contribution lisse et de
longue portée.

La troisième image correspond à la méthode développée pendant le stage.
La composante basse fréquence est calculée à résolution spatiale et
angulaire réduite, puis la composante haute fréquence locale est ajoutée
séparément. Cette recombinaison constitue la reconstruction finale obtenue
avec la nouvelle stratégie.


\section{Comparaison avec la correction complète de référence}

La comparaison la plus importante est celle entre la nouvelle méthode
multirésolution et la correction complète utilisée comme référence.

% ============================================================
% FIGURE 8
% ============================================================

% Deux panneaux :
% (a) correction complète de référence
% (b) nouvelle correction LF + HF
%
% Même coupe, même crop, même grayscale.
%
% Nom conseillé :
% figures/08_reconstruction_reference_comparison.png

\begin{figure}[H]
\centering

\begin{tikzpicture}[
    node distance=1.8cm
]

% Image de référence
\node (img1) {
    \includegraphics[width=0.32\textwidth]
    {figures/ref.png}
};

% Nouvelle méthode
\node[right=1.8cm of img1] (img2) {
    \includegraphics[width=0.32\textwidth]
    {figures/correc new.png}
};

% Labels sous les images
\node[below=0.2cm of img1, align=center]
{\small Reference\\full-resolution correction};

\node[below=0.2cm of img2, align=center]
{\small New LF + HF\\correction};

\end{tikzpicture}

\caption{
Comparaison entre la reconstruction obtenue avec la correction complète de
référence et celle obtenue avec la nouvelle stratégie LF + HF. Les deux
images correspondent à la même coupe reconstruite, avec le même recadrage
et la même plage de niveaux de gris, ce qui permet une comparaison directe.
}

\label{fig:reconstruction_reference}

\end{figure}



L'objectif de la nouvelle méthode n'est pas de produire une correction
différente, mais de retrouver aussi fidèlement que possible l'effet de la
correction complète tout en évitant son coût de calcul et de stockage à
pleine résolution.

Visuellement, la reconstruction obtenue avec la stratégie LF + HF reste
très proche de celle obtenue avec la correction complète de référence. Les
principales structures de la coupe sont conservées et aucune dégradation
évidente n'apparaît dans cette comparaison.

Cette comparaison doit cependant être interprétée comme une validation
visuelle. Les métriques quantitatives présentées dans les sections
précédentes ont été calculées sur les composantes de correction et non
directement entre ces deux images reconstruites.

Une fois vérifié que la nouvelle stratégie reste visuellement cohérente
avec la correction de référence, la dernière question concerne donc le gain
obtenu : quelle réduction apporte-t-elle réellement en termes d'espace
disque et de temps de traitement ?




% ============================================================
% CHAPITRE 12
% ============================================================

\chapter{Performances mesurées}

Après avoir vérifié que la nouvelle correction reste visuellement cohérente
avec la correction complète de référence, il reste à évaluer le gain
pratique apporté par cette nouvelle organisation du pipeline.

Les mesures finales ont été réalisées pour la configuration étudiée sur un
nœud de calcul de 96 cœurs. Elles concernent trois aspects importants :
l'espace disque nécessaire, le temps de création du digest et le temps de
reconstruction.

\begin{table}[H]
\centering

\caption{Performances mesurées pour la méthode de référence et la nouvelle
stratégie multirésolution.}

\label{tab:performance}

\begin{tabular}{lccc}
\toprule

\textbf{Grandeur}
&
\textbf{Référence}
&
\textbf{Nouvelle méthode}
&
\textbf{Gain}
\\

\midrule

Empreinte disque
&
4.7 To
&
80 Go
&
$\approx 59\times$ plus faible
\\

Création du digest
&
32 min
&
4 min
&
$8\times$ plus rapide
\\

Reconstruction
&
118 min
&
93 min
&
$\approx 21.2\,\%$ plus rapide
\\

\bottomrule
\end{tabular}

\end{table}

Le gain le plus important concerne l'empreinte disque. Le volume de données
intermédiaires passe d'environ 4.7 To à 80 Go, soit une réduction d'un
facteur proche de 59.

Le temps nécessaire à la création du digest passe quant à lui de
32 minutes à 4 minutes. Cette étape est donc environ huit fois plus rapide
avec la nouvelle organisation.

Enfin, le temps de reconstruction complet diminue de 118 minutes à
93 minutes, soit un gain de 25 minutes et une réduction d'environ
21.2\,\%.

Ces résultats montrent que la stratégie multirésolution ne constitue pas
uniquement une réorganisation numérique de la correction. La séparation
entre une composante basse fréquence réduite et une composante haute
fréquence locale entraîne également une diminution concrète du coût du
pipeline, aussi bien en termes de stockage que de temps de traitement.

Ces valeurs correspondent à la configuration, au jeu de données et au
matériel utilisés pour ces mesures. Elles ne doivent donc pas être
interprétées comme des performances universelles.


% ============================================================
% CHAPITRE 13
% ============================================================

\chapter{Outils utilisés et compétences développées}

\section{Environnement et outils utilisés}

Le stage m'a amenée à utiliser un ensemble assez large d'outils, allant du
calcul scientifique à l'édition de code, la visualisation d'images, la gestion
de fichiers et l'intégration dans la chaîne de reconstruction de l'ESRF.

\subsection{Programmation et calcul scientifique}

\begin{itemize}
    \item \textbf{Python} pour le prototypage, les tests et le traitement
    basse fréquence ;
    
    \item \textbf{NumPy} pour la manipulation des tableaux et des données ;
    
    \item \textbf{SciPy} pour le filtrage, l'interpolation et certaines
    opérations numériques ;
    
    \item \textbf{h5py} pour la lecture et l'écriture des fichiers HDF5 ;
    
    \item \textbf{pyFFTW} et \textbf{FFTW} pour les convolutions par
    transformée de Fourier ;
    
    \item \textbf{Matplotlib} pour les figures, profils et validations ;
    
    \item \textbf{multiprocessing} pour certains calculs parallélisés ;
    
    \item \textbf{C++} pour l'intégration de la composante haute fréquence
    dans \texttt{nabuxx} ;
    
    \item \textbf{xtensor} pour la manipulation des tableaux numériques
    côté C++ ;
    
    \item \textbf{HighFive} pour l'accès aux fichiers HDF5 côté C++.
\end{itemize}

\subsection{Visualisation et traitement d'images}

\begin{itemize}
    \item \textbf{ImageJ} pour l'inspection et la comparaison des images et
    reconstructions ;
    
    \item \textbf{silx} pour la visualisation de données scientifiques et
    de fichiers HDF5 ;
    
    \item analyse de fichiers \texttt{.tif}, profils d'intensité,
    comparaison de régions d'intérêt et contrôle visuel des reconstructions.
\end{itemize}

\subsection{Édition et comparaison de fichiers}

\begin{itemize}
    \item \textbf{Emacs} et \textbf{nano} pour l'édition de scripts et de
    fichiers directement sur les machines de calcul ;
    
    \item \textbf{Meld} pour comparer différentes versions de fichiers et
    identifier précisément les modifications introduites dans le code ;
    
    \item navigation et modification de fichiers en environnement Linux.
\end{itemize}

\subsection{Environnements de calcul}

\begin{itemize}
    \item utilisation d'environnements virtuels Python afin d'isoler les
    dépendances nécessaires aux différents scripts ;
    
    \item installation et gestion de bibliothèques scientifiques ;
    
    \item travail sur des machines Linux distantes ;
    
    \item lancement et suivi de calculs depuis le terminal ;
    
    \item gestion de données expérimentales volumineuses et de fichiers
    intermédiaires.
\end{itemize}

\subsection{Terminal Linux}

Au cours du stage, j'ai progressivement utilisé le terminal comme un outil
de travail quotidien. Parmi les commandes fréquemment utilisées :

\begin{verbatim}
pwd
ls
cd
cp
mv
rm
mkdir
grep
find
head
tail
less
ps
\end{verbatim}

Ces commandes ont été utilisées pour naviguer dans les répertoires,
rechercher des fichiers ou des portions de code, inspecter des résultats,
organiser les données et suivre l'exécution des programmes.

\subsection{Versionnement et gestion du code}

\begin{itemize}
    \item utilisation de \textbf{Git} pour suivre les modifications ;
    
    \item commandes telles que \texttt{git status}, \texttt{git diff},
    \texttt{git add} et \texttt{git commit} ;
    
    \item utilisation de \textbf{GitHub} pour centraliser les scripts,
    figures et documents partageables liés au projet.
\end{itemize}



\section{Compétences développées}

\subsection{Relier physique et calcul numérique}

L'un des principaux acquis du stage a été d'apprendre à partir d'un
phénomène physique réel pour construire progressivement une stratégie
numérique.

Le travail a commencé par la diffusion optique dans le système de détection,
puis par sa représentation sous forme de kernel et de convolution. L'étude
des différentes échelles spatiales a ensuite conduit à la séparation basse
fréquence / haute fréquence, puis à deux traitements numériques adaptés aux
propriétés de chacune de ces composantes.

L'optimisation n'a donc pas été introduite comme une simple astuce de
programmation : elle découle directement de la compréhension de la
structure physique et spatiale de la correction.

\subsection{Construire et valider une méthode progressivement}

Le développement a également demandé une démarche de validation
progressive. Avant de modifier la chaîne complète, chaque nouvelle idée a
été testée sur des cas contrôlés : binning, interpolation entre projections,
séparation fréquentielle, taille de la marge haute fréquence puis
recombinaison des deux branches.

Cette méthode m'a appris à ne pas juger une optimisation uniquement sur son
temps d'exécution, mais à vérifier d'abord l'erreur qu'elle introduit et à
identifier précisément la composante qui devient limitante.

\subsection{Manipuler des données expérimentales volumineuses}

Le projet m'a confrontée à des volumes de données beaucoup plus importants
que dans mes projets universitaires précédents. Cela m'a amenée à prendre
en compte non seulement le temps de calcul, mais aussi l'espace disque, les
accès aux fichiers, les allocations mémoire, la réutilisation des kernels
et le coût des données intermédiaires.

Cette expérience a donné une dimension plus concrète à la notion de
performance en calcul scientifique.

\subsection{Intégrer une solution dans un environnement de production}

Une autre compétence importante a été de passer du prototype à
l'intégration.

Une solution qui fonctionne dans un script isolé ne suffit pas lorsqu'elle
doit être insérée dans une chaîne utilisée pour traiter de véritables
acquisitions. Il faut gérer les formats de fichiers, les paramètres, les
dimensions, les blocs de projections, les interfaces Python/C++ et la
transmissibilité du code.

Cette étape a constitué ma première expérience significative d'intégration
dans une base C++ scientifique existante.

\subsection{Communication scientifique}

Enfin, ce stage m'a également demandé de synthétiser un travail mêlant
physique, traitement d'image et informatique scientifique afin de le rendre
compréhensible à différents interlocuteurs.

La préparation du rapport et de la présentation finale m'a notamment
obligée à distinguer les éléments essentiels du projet des détails
d'implémentation, et à construire une explication allant du phénomène
physique jusqu'aux gains mesurés.




\chapter*{Conclusion}
\addcontentsline{toc}{chapter}{Conclusion}

Ce stage à l'ESRF m'a permis de découvrir concrètement la manière dont un
problème physique peut être transformé en problème numérique, puis intégré
dans une chaîne de traitement utilisée sur des données expérimentales
réelles.

Le travail réalisé m'a amenée à combiner plusieurs aspects que j'avais
jusqu'alors abordés séparément : compréhension physique du système de
détection, traitement d'image, méthodes numériques, calcul scientifique,
gestion de grands volumes de données et développement dans une base de code
existante.

J'ai également acquis une expérience beaucoup plus pratique des
environnements de recherche informatique : Linux, travail sur machines
distantes, Python, C++, FFTW, HDF5, Git, outils de visualisation et
débogage. Une partie importante du stage a consisté à apprendre à tester une
idée progressivement, à identifier précisément l'origine d'une erreur et à
vérifier qu'une optimisation restait cohérente avec le résultat physique
attendu.

Ce stage s'inscrit dans la continuité de mon parcours scientifique. Après
une première expérience à l'IRAM consacrée à l'analyse de données
interférométriques ALMA, l'ESRF m'a permis de découvrir un autre grand
instrument scientifique et de travailler davantage à l'interface entre
physique, instrumentation et calcul.

À la suite de cette expérience, j'intègre le Master Nanophysics and Quantum
Physics de l'Université Grenoble Alpes ainsi que le Quantum Program de la
Graduate School. Je souhaite poursuivre dans des projets combinant
physique, instrumentation, traitement numérique et technologies quantiques,
en m'appuyant sur les compétences développées progressivement au cours de
mes expériences à l'IRAM puis à l'ESRF.


\appendix

\chapter{Fichiers, scripts et dépôt du projet}

Les principaux fichiers de production utilisés au cours du stage sont :

\begin{itemize}

    \item \texttt{diffusion\_correction.py} ;
    \item \texttt{diffusion\_parameters.json} ;
    \item \texttt{diffusion\_correction.h5} ;
    \item \texttt{kernel.h5} ;
    \item \texttt{etfscan.cpp} ;
    \item \texttt{etfscan.hpp}.

\end{itemize}

Plusieurs scripts indépendants ont également été développés pendant la
phase de validation afin d'étudier successivement le binning, les
interpolations, la séparation basse fréquence / haute fréquence, la
correction locale et la recombinaison multirésolution.

Le script :

\begin{center}
\texttt{23\_full\_correction\_spatial4\_angular4\_local\_high.py}
\end{center}

constitue notamment une synthèse de la phase de validation du pipeline.

L'ensemble des éléments pouvant être partagés, notamment les scripts de
validation, les figures et les documents de présentation du projet, est
regroupé dans le dépôt GitHub associé au stage.

\begin{center}
\textbf{Dépôt GitHub : }\\
\url{ https://github.com/alouaniranou2004-droid}
\end{center}

\end{document}
